\kchapter{Resources and Scheduling}{quota}

\begin{objectives}
\begin{itemize}
  \item CPU and memory \emph{requests} and \emph{limits}, and their units;
  \item Quality-of-Service classes and what happens under pressure;
  \item LimitRanges and ResourceQuotas for namespaces;
  \item controlling placement: nodeSelector, affinity, taints and tolerations,
        topology spread, priority and preemption;
  \item PodDisruptionBudgets and node draining.
\end{itemize}
\end{objectives}

\section{Requests and limits}

Each container can declare two numbers per resource:

\begin{itemize}
  \item \kterm{request} --- what the container is \emph{guaranteed}. The
        scheduler only places a Pod on a node with enough unrequested capacity.
  \item \kterm{limit} --- the \emph{maximum} it may use.
\end{itemize}

\begin{yamlcode}{resources.yaml}
resources:
  requests:
    cpu: 250m          # 250 millicores = 1/4 of a CPU core
    memory: 256Mi      # mebibytes (Mi, Gi) -- not MB
  limits:
    cpu: "1"           # optional, see the Pro Tip below
    memory: 512Mi
\end{yamlcode}

\begin{center}
\ktablestart
\begin{tabularx}{\linewidth}{P{2.0cm} L L}
\kheader \kth{} & \kth{CPU (compressible)} & \kth{Memory (incompressible)} \\
Units & \kc{1} = one core; \kc{500m} = half a core. & \kc{Mi}, \kc{Gi} (powers of 2); \kc{M}, \kc{G} (powers of 10). \\
Over limit & The container is \textbf{throttled} (slowed down). & The container is \textbf{OOM-killed} and restarted. \\
\end{tabularx}
\end{center}

\begin{protip}
Always set \textbf{memory requests and limits} (often equal). Always set
\textbf{CPU requests}. Many teams \emph{omit CPU limits} for latency-sensitive
services, because CPU throttling causes hard-to-diagnose slowness while idle
cores sit unused. Measure first with \kc{kubectl top} and your monitoring.
\end{protip}

\section{Quality of Service classes}

From requests and limits, Kubernetes assigns each Pod a QoS class. When a node
runs out of memory, the kubelet evicts Pods in this order: BestEffort first,
then Burstable, and Guaranteed last.

\begin{center}
\ktablestart
\begin{tabularx}{\linewidth}{P{2.6cm} L P{3.0cm}}
\kheader \kth{Class} & \kth{Rule} & \kth{Eviction priority} \\
\textcolor{k8sgreen}{\textbf{Guaranteed}} & Every container has CPU and memory requests \emph{equal to} limits. & Last to go. \\
\textcolor{k8sorange}{\textbf{Burstable}} & At least one request or limit set, but not Guaranteed. & Middle. \\
\textcolor{k8sred}{\textbf{BestEffort}} & No requests or limits at all. & First to go. \\
\end{tabularx}
\end{center}

\section{Namespace guard rails}

\subsection{\kicon{limits}\ LimitRange: defaults per container}

\begin{yamlcode}{limitrange.yaml}
apiVersion: v1
kind: LimitRange
metadata:
  name: defaults
  namespace: shop
spec:
  limits:
    - type: Container
      defaultRequest: { cpu: 100m, memory: 128Mi }   # applied if omitted
      default:        { memory: 256Mi }              # default limit
      max:            { cpu: "2", memory: 2Gi }
\end{yamlcode}

\subsection{\kicon{quota}\ ResourceQuota: a budget per namespace}

\begin{yamlcode}{quota.yaml}
apiVersion: v1
kind: ResourceQuota
metadata:
  name: shop-budget
  namespace: shop
spec:
  hard:
    requests.cpu: "8"
    requests.memory: 16Gi
    limits.memory: 32Gi
    pods: "60"
    persistentvolumeclaims: "10"
    services.loadbalancers: "1"
\end{yamlcode}

\begin{bashcode}{terminal}
kubectl describe quota shop-budget -n shop     # used vs hard
\end{bashcode}

\section{Controlling where Pods run}

\subsection{nodeSelector: the simple way}

\begin{bashcode}{terminal}
kubectl label node lab-worker disktype=ssd
\end{bashcode}

\begin{yamlcode}{nodeselector.yaml}
spec:
  nodeSelector:
    disktype: ssd
\end{yamlcode}

\subsection{Node affinity: the expressive way}

\begin{yamlcode}{node-affinity.yaml}
spec:
  affinity:
    nodeAffinity:
      requiredDuringSchedulingIgnoredDuringExecution:     # hard rule
        nodeSelectorTerms:
          - matchExpressions:
              - key: kubernetes.io/arch
                operator: In
                values: [amd64, arm64]
      preferredDuringSchedulingIgnoredDuringExecution:    # soft preference
        - weight: 80
          preference:
            matchExpressions:
              - key: topology.kubernetes.io/zone
                operator: In
                values: [eu-west-1a]
\end{yamlcode}

\subsection{Pod affinity and anti-affinity}

Place Pods \emph{relative to other Pods}: keep a cache near its app, or keep
replicas of the same app apart.

\begin{yamlcode}{pod-anti-affinity.yaml}
spec:
  affinity:
    podAntiAffinity:
      preferredDuringSchedulingIgnoredDuringExecution:
        - weight: 100
          podAffinityTerm:
            topologyKey: kubernetes.io/hostname    # "not on the same node"
            labelSelector:
              matchLabels: { app: api }
\end{yamlcode}

\subsection{Topology spread constraints}

The modern, more precise way to spread replicas across zones and nodes:

\begin{yamlcode}{topology-spread.yaml}
spec:
  topologySpreadConstraints:
    - maxSkew: 1
      topologyKey: topology.kubernetes.io/zone
      whenUnsatisfiable: ScheduleAnyway
      labelSelector:
        matchLabels: { app: api }
    - maxSkew: 1
      topologyKey: kubernetes.io/hostname
      whenUnsatisfiable: DoNotSchedule
      labelSelector:
        matchLabels: { app: api }
\end{yamlcode}

\subsection{Taints and tolerations}

Affinity \emph{attracts} Pods to nodes. \kterm{Taints} \emph{repel} them: a
tainted node accepts only Pods that \kterm{tolerate} the taint.

\begin{analogy}
A taint is a ``Staff only'' sign on a door. Most people (Pods) stay out. A
toleration is the staff badge. Note that the badge \emph{allows} entry but does
not \emph{force} it --- combine with affinity if you want the Pod to go there.
\end{analogy}

\begin{bashcode}{terminal}
kubectl taint nodes gpu-node-1 gpu=true:NoSchedule
kubectl taint nodes gpu-node-1 gpu=true:NoSchedule-     # remove (trailing dash)
\end{bashcode}

\begin{yamlcode}{toleration.yaml}
spec:
  tolerations:
    - key: gpu
      operator: Equal
      value: "true"
      effect: NoSchedule
  nodeSelector:
    gpu: "true"
\end{yamlcode}

Effects: \kc{NoSchedule} (do not place new Pods), \kc{PreferNoSchedule} (try
not to), \kc{NoExecute} (also evict running Pods that do not tolerate it).

\subsection{Priority and preemption}

\begin{yamlcode}{priorityclass.yaml}
apiVersion: scheduling.k8s.io/v1
kind: PriorityClass
metadata:
  name: business-critical
value: 100000
globalDefault: false
description: "Customer-facing services; may preempt batch work."
---
# in a Pod spec:
# priorityClassName: business-critical
\end{yamlcode}

If a high-priority Pod cannot be scheduled, the scheduler may \emph{evict}
lower-priority Pods to make room.

\section{Voluntary disruptions}

Node upgrades and \kc{kubectl drain} evict Pods on purpose. A
\kterm{PodDisruptionBudget} (PDB) tells Kubernetes how many Pods of an app must
stay up during such voluntary disruptions.

\begin{yamlcode}{pdb.yaml}
apiVersion: policy/v1
kind: PodDisruptionBudget
metadata:
  name: api-pdb
spec:
  minAvailable: 2           # or maxUnavailable: 1
  selector:
    matchLabels: { app: api }
\end{yamlcode}

\begin{bashcode}{terminal}
kubectl cordon lab-worker                              # no new Pods here
kubectl drain lab-worker --ignore-daemonsets --delete-emptydir-data
kubectl uncordon lab-worker
\end{bashcode}

\begin{recap}
\begin{itemize}
  \item Requests drive scheduling; limits cap usage. Memory over limit = OOMKill.
  \item QoS classes decide who gets evicted first.
  \item LimitRanges set defaults; ResourceQuotas cap a namespace.
  \item nodeSelector/affinity attract, taints repel, topology spread distributes.
  \item PDBs keep apps available during drains and upgrades.
\end{itemize}
\end{recap}
